\documentclass[submission,copyright,creativecommons]{eptcs}
\providecommand{\event}{ACT 2026}

\usepackage{iftex}

\ifpdf
  \usepackage{underscore}
  \usepackage[T1]{fontenc}
\fi

\usepackage{amsmath}
\usepackage{amssymb}
\usepackage{amsthm}

\title{From Games to Graphs:\\Categorical Composition of Genetic Algorithms Across Domains}

\author{Robin Langer
\institute{Universit\'{e} Paris Est -- Marne la Vall\'{e}e}
\email{langer.robin@gmail.com}
\and
Lyra Vega
\institute{Independent Researcher}
\and
Nick Meinhold
\institute{Independent Researcher}
\and
Cale Gibbard
\institute{Independent Researcher}
}

\newcommand{\titlerunning}{From Games to Graphs}
\newcommand{\authorrunning}{R. Langer, L. Vega, N. Meinhold \& C. Gibbard}

\hypersetup{
  bookmarksnumbered,
  pdftitle    = {\titlerunning},
  pdfauthor   = {\authorrunning},
  pdfsubject  = {ACT 2026},
  pdfkeywords = {category theory, genetic algorithms, Kleisli morphisms,
                 evolutionary computation, compositional optimization}
}

\newtheorem{conjecture}{Conjecture}

\begin{document}
\maketitle

\begin{abstract}
Categorical formalizations of neural networks, reinforcement learning,
and compositional RL have shown that how optimization components compose
determines system-level behavior. Evolutionary computation---the last
major paradigm---has remained outside this categorical landscape. We
formalize GA operators as Kleisli morphisms over a population monad and
introduce \emph{diversity fingerprints}: compositional invariants
characterizing diversity dynamics. The same composition produces the same
fingerprint across two unrelated domains (Cohen's $d = 4.34$,
$p < 3.69 \times 10^{-11}$). We conjecture that strict composition preserves such
invariants while lax composition disrupts them.
\end{abstract}

\section{The Optimization Zoo}

Over the past three years, three independent programs have applied
category theory to optimization, arriving at a common structural insight:
composition determines behavior.
Gavranovi\'{c} et al.~\cite{gavranovic2024categorical} showed that deep
learning architectures are constructions in $\mathbf{Para}$, where
backpropagation is Kleisli composition.
Hedges and Sakamoto~\cite{hedges2024reinforcement} formalized RL within
categorical cybernetics, expressing the Bellman equation as a natural
transformation between functors.
Bakirtzis et al.~\cite{bakirtzis2024compositional} extended this to
multi-agent RL, deriving safety guarantees from categorical composition
of MDPs.

One major optimization paradigm remains outside this landscape:
\emph{evolutionary computation}. Genetic algorithms have been studied for
over five decades~\cite{holland1975adaptation,goldberg1989genetic}, yet
their formalization has stayed set-theoretic. The compositional structure
of a GA pipeline---selection, crossover, mutation---has been treated as
an implementation detail. This talk presents our work filling that gap.

\section{GA Operators as Kleisli Morphisms}

Following Moggi~\cite{moggi1991notions}, we define the \emph{evolution
monad} $T = \mathsf{Reader}(\Gamma) \times \mathsf{Writer}(\Lambda)
\times \mathsf{State}(S)$, where $\Gamma$ is a configuration context,
$\Lambda$ a log monoid of diversity metrics, and $S$ the state of a
pseudorandom generator. A GA operator is a Kleisli morphism
$f : \mathsf{Pop}(A) \to T(\mathsf{Pop}(B))$, and the Kleisli category
$\mathbf{Kl}(T)$ supports composition at three levels: operators into
pipelines, pipelines into strategies, and strategies into
multi-population systems via an island functor $\mathcal{I}$.

The framework is implemented in Haskell, where the monad stack is
\texttt{ReaderT GAConfig (StateT StdGen (Writer GALog))}. The same
strategy code operates unchanged on checkers weight vectors
(\texttt{[Double]}) and maze topologies (\texttt{[Bool]})---only the
fitness function changes. Because a pipeline is itself a Kleisli
morphism, it plugs into any strategy combinator without modification.

\section{Key Result: Diversity Fingerprints}

We tested four composition strategies on two unrelated domains---checkers
evaluation and maze topology evolution. Each strategy applies the same
operators in a different compositional arrangement. The central finding
is that each composition produces a characteristic \emph{diversity
fingerprint}---a temporal pattern in population diversity invariant
across domains. We observe four distinct fingerprints: monotonic decline,
spike-crash-rebound, stable oscillation, and spike-collapse.

Statistical validation over 30 seeds confirms cross-domain stability:
Cohen's $d = 4.34$ (effect size: the strict and lax distributions are
4+ standard deviations apart---a massive separation), with
$p < 3.69 \times 10^{-11}$ (less than 1 in 27 billion chance this is
random). Compositional structure is the dominant predictor of diversity
dynamics, not domain-specific properties.

\section{Conjecture and Connections}

\begin{conjecture}[Compositional invariance]
For any monoidal category of optimization components, strict composition
preserves compositional invariants; lax composition disrupts them.
\end{conjecture}

\noindent
Evidence comes from all four zoo entries. In our framework, the island
functor $\mathcal{I}$ preserves fingerprints when migration is absent
(strict) and disrupts them when migration introduces cross-island
information flow (lax). This mirrors Ghani et al.'s approximate game
theory~\cite{ghani2025approximate}, where exact Nash equilibria compose
but $\varepsilon$-approximate equilibria do not---a strict/lax
distinction in a different categorical setting.

\paragraph{Talk outline.}
The talk will present: (1)~the optimization zoo and why evolutionary
computation is the missing piece; (2)~the Kleisli formalization of GA
operators; (3)~the diversity fingerprint result with statistical
evidence; and (4)~the strict/lax conjecture and its connections to
approximate game theory. Our work completes a categorical picture of four
major optimization paradigms. We believe this convergence---where
mathematical structure makes precise predictions confirmed by
computation---is a natural fit for ACT.

\bibliographystyle{eptcs}
\bibliography{references}

\end{document}
